% Guerra arancelaria (Johnson, 1953): estado final. Curvas de oferta neta con
% la notación del 3.B.5: X la exporta A (la importa B), Y la exporta B.
% Caso simétrico (para que se vea la idea, no un caso general):
%   ON_B: Y = g(X), ON_A: X = g(Y), con g(t) = 6t/(1+t) - 0,15t (cóncava,
%         con tramo inelástico al final)
%   E_0: g(t) = t -> t = 4,217; RRI_0 = p_X/p_Y = 1 (la bisectriz)
%   Con aranceles, las 2 curvas se contraen: ON_B': Y = 0,6 g(X);
%         ON_A': X = 0,6 g(Y)
%   E_N: 0,6 g(t) = t -> t = 2,303 (menos comercio, pero no nulo)
%   Pendiente de ON_B' en E_N: 0,6 (6/(1+t)^2 - 0,15) = 0,240
%   CIC_A: Y = Y_N + 0,24 (X - X_N) + 0,13 (X - X_N)^2, convexa, tangente a
%          ON_B' en E_N y por encima de ella (A prefiere más Y y menos X)
%   CIC_B: X = X_N + 0,24 (Y - Y_N) + 0,13 (Y - Y_N)^2, simétrica, tangente a
%          ON_A' en E_N
%   Al ser simétrico, la relación real de intercambio no cambia y los 2
%   países pierden frente a E_0.
% Pasos (rehacen las Imágenes 4-6 del tema anterior): johnson-paso1.tex
%   (\tceepaso=2: libre comercio, con las CIC de A y B tangentes en E_0),
%   johnson-paso2.tex (\tceepaso=3: arancel óptimo de A, E_1 sobre ON_B,
%   E_1 = (2,515; 3,916), RRI_1 = 1,557, pendiente de ON_B en E_1 = 0,336) y
%   este fichero completo (equilibrio de Nash). \solopaso{k}{...} dibuja algo
%   solo hasta el paso k, para que las CIC de pasos anteriores no estorben.
% Autor: Víctor Gutiérrez Marcos
\documentclass[tikz,border=4pt]{standalone}
% ==========================================================================
% tikz-tcee.tex — Estilos comunes de los gráficos TikZ del temario
% Autor: Víctor Gutiérrez Marcos
%
% Cada gráfico es un fichero standalone en fuentes/<tema>/graficos/:
%
%   \documentclass[tikz,border=4pt]{standalone}
%   % ==========================================================================
% tikz-tcee.tex — Estilos comunes de los gráficos TikZ del temario
% Autor: Víctor Gutiérrez Marcos
%
% Cada gráfico es un fichero standalone en fuentes/<tema>/graficos/:
%
%   \documentclass[tikz,border=4pt]{standalone}
%   % ==========================================================================
% tikz-tcee.tex — Estilos comunes de los gráficos TikZ del temario
% Autor: Víctor Gutiérrez Marcos
%
% Cada gráfico es un fichero standalone en fuentes/<tema>/graficos/:
%
%   \documentclass[tikz,border=4pt]{standalone}
%   \input{../../../latex/tikz-tcee}
%   \begin{document}
%   \begin{tikzpicture}[tcee]
%     \ejes{Q}{P}{6}{5}
%     \draw[curva1] (0.5,4.5) -- (5,0.8) node[etiqueta,right] {$D$};
%     \capa{2}{\draw[curva2] ...;}   % en el vídeo aparece en el paso 2
%   \end{tikzpicture}
%   \end{document}
%
% Paleta: la de la web (morado, dorado, salvia) más tres colores de apoyo
% distinguibles en blanco y negro por el trazo.
% ==========================================================================
\usepackage[T1]{fontenc}
\usepackage{newpxtext,newpxmath}
\usepackage{xcolor}
\usepackage{tikz,pgfplots}
\pgfplotsset{compat=1.18}
\usetikzlibrary{arrows.meta,calc,positioning,patterns,decorations.pathreplacing,intersections,shapes.misc,backgrounds}

\definecolor{tceemorado}{HTML}{5F2987}
\definecolor{tceedorado}{HTML}{B8860B}
\definecolor{tceeverde}{HTML}{3C7D3A}
\definecolor{tceeazul}{HTML}{1F5F99}
\definecolor{tceerojo}{HTML}{B22222}
\definecolor{tceenaranja}{HTML}{E07B28}
\definecolor{tceegris}{HTML}{767171}
\definecolor{tceesalvia}{HTML}{E2EFD9}

% Capas para el vídeo: \capa{n}{…} dibuja su contenido solo a partir del paso n.
% En el PDF, la web y Word se ve todo; generar-video.py compila el gráfico una
% vez por paso con \def\tceepaso{k} para que cada elemento aparezca al narrarlo.
% Vectores con flecha encima, como en el texto (tcee.sty): \vect{x}, \vect{p}
\providecommand{\vect}[1]{\vec{#1}}
\providecommand{\tceepaso}{1000}
\newcommand{\capa}[2]{\ifnum#1>\tceepaso\relax\else#2\fi}
% Gráficos por pasos del mismo tamaño (propuesta del 3.A.29, 11/11/2026): con
% \def\tceefantasma{1} antes del \input (p. ej., x-paso1.tex:
% «\def\tceepaso{3}\def\tceefantasma{1}\input{x.tex}»), las capas posteriores
% al paso se dibujan invisibles y ocupan su sitio. Sin \tceefantasma (el vídeo),
% no cambia nada.
\ifdefined\tceefantasma
  \renewcommand{\capa}[2]{\ifnum#1>\tceepaso\relax\begin{scope}[opacity=0]#2\end{scope}\else#2\fi}
\fi

\tikzset{
  tcee/.style={line cap=round,line join=round,font=\small,>={Stealth[length=2.2mm]}},
  eje/.style={->,thick,black},
  curva1/.style={very thick,tceemorado},
  curva2/.style={very thick,tceeazul},
  curva3/.style={very thick,tceeverde},
  curva4/.style={very thick,tceedorado},
  curvanueva/.style={very thick,tceerojo},
  desplazada/.style={very thick,dashed},
  guia/.style={thin,dashed,tceegris},
  punto/.style={circle,fill=black,inner sep=1.3pt},
  etiqueta/.style={font=\small,inner sep=2pt},
  flecha/.style={->,thick,tceerojo},
  area/.style={fill=tceemorado,fill opacity=0.18,draw=none},
  area2/.style={fill=tceedorado,fill opacity=0.22,draw=none},
}

% \ejes{etiqueta x}{etiqueta y}{largo x}{alto y}
\newcommand{\ejes}[4]{%
  \draw[eje] (0,0) -- (#3,0) node[below left=2pt and 0pt] {#1};
  \draw[eje] (0,0) -- (0,#4) node[left] {#2};
  \node[below left] at (0,0) {$0$};}
% \proyeccion{(x,y)}{etq x}{etq y}: líneas guía a los ejes con etiquetas
\newcommand{\proyeccion}[3]{%
  \draw[guia] let \p1=#1 in (\x1,0) node[below] {#2} |- (0,\y1) node[left] {#3};}

% --------------------------------------------------------------------------
% Paleta «dictamen» (dic…): gráficos de datos de la Guía del dictamen (1D01) y
% de cualquier gráfico de coyuntura que quiera los mismos colores. Un concepto,
% un color en todos los gráficos (CRITERIOS §7, «tres reglas», 24/10/2026);
% tabla completa en fuentes/1D01/_trabajo/paleta-1D01.md. No cambia nada de
% lo anterior: son colores y estilos nuevos.
% Actividad: el PIB (y sus variantes: PIB por persona, PIB mundial), morado.
\definecolor{dicpib}{HTML}{5F2987}
% Demanda interna, gama cálida: consumo (naranjas: privado oscuro, público
% claro), inversión (verde) y existencias (gris con trama, aparte).
\definecolor{dicconspriv}{HTML}{C76A1F}
\definecolor{dicconspub}{HTML}{F5C690}
\definecolor{dicfbcf}{HTML}{2A6B3C}
\definecolor{dicexistencias}{HTML}{B4B4B4}
% Demanda externa y comercio, gama fría (azul claro, que se separa del morado)
\definecolor{dicexterna}{HTML}{6AA6DA}
% Zonas: zona del euro (azul marino), España (rojo), resto de países (grises)
\definecolor{dicze}{HTML}{17406D}
\definecolor{dicespana}{HTML}{C8382E}
\definecolor{dicotro}{HTML}{8A8A8A}
\definecolor{dicotroclaro}{HTML}{B5B5B5}
\definecolor{dicotromasclaro}{HTML}{DADADA}
% Precios: inflación general (negro, continuo) y subyacente (negro,
% discontinuo); componentes: energía (bermellón), alimentos (amarillo) y los
% de la subyacente en azules (bienes industriales claro, servicios oscuro).
\definecolor{dicinflacion}{HTML}{1A1A1A}
\definecolor{dicenergia}{HTML}{D55E00}
\definecolor{dicalimentos}{HTML}{F2D45C}
\definecolor{dicbine}{HTML}{7FB2E0}
\definecolor{dicservicios}{HTML}{1F4E89}
% Trabajo: empleo (verde azulado), horas (el mismo, claro), tasa de actividad
% (verde azulado oscuro), paro (teja); productividad (dorado de la web).
\definecolor{dicempleo}{HTML}{1B7F6B}
\definecolor{dichoras}{HTML}{86CBB8}
\definecolor{dicactividad}{HTML}{0F5247}
\definecolor{dicparo}{HTML}{A23B2A}
\definecolor{dicproductividad}{HTML}{B8860B}
% Hogares: renta disponible (pizarra) y tasa de ahorro (negro); el consumo es
% dicconspriv.
\definecolor{dicrenta}{HTML}{4A5A6A}
\definecolor{dicahorro}{HTML}{1A1A1A}
% Política monetaria: tipos del BCE en azules (depósito, el de la zona del
% euro; operaciones principales, medio; facilidad marginal, claro); balance:
% préstamos a los bancos (dorado claro), resto del activo (gris claro).
\definecolor{dicdeposito}{HTML}{17406D}
\definecolor{dicopf}{HTML}{4F86C6}
\definecolor{dicfmc}{HTML}{9EC5EA}
\definecolor{dicprestamos}{HTML}{D9A520}
\definecolor{dicresto}{HTML}{CFCFCF}
% Ciclo: series adelantadas, coincidentes (el PIB) y retrasadas
\definecolor{dicadelantada}{HTML}{D55E00}
\definecolor{dicretrasada}{HTML}{6AA6DA}
% Referencias: rejilla, previsión, sombreados
\definecolor{dicreferencia}{HTML}{767171}

\tikzset{
  % barra de un componente: \addplot[dicbarra=dicfbcf]
  dicbarra/.style={fill=#1, draw=#1!65!black, line width=0.3pt},
  % existencias: gris con trama, para que se distinga también en gris
  dicbarratrama/.style={fill=#1, draw=#1!65!black, line width=0.3pt,
    postaction={pattern=north east lines, pattern color=white}},
  % línea de una serie; con halo blanco cuando pasa sobre barras
  diclinea/.style={very thick, #1},
  diclineahalo/.style={very thick, #1, preaction={draw=white, line width=3.4pt, solid}},
  % puntos de una serie (capa aparte, después de la línea)
  dicpuntos/.style={only marks, mark=*, mark size=1.8pt,
    mark options={fill=white, draw=#1, line width=0.9pt, solid}},
  % línea del cero o de un umbral, después de las barras
  diccero/.style={thin, black},
  dicumbral/.style={thin, dashed, dicreferencia},
}
% Capas de los gráficos de datos (CRITERIOS §7, «tres reglas»), de abajo
% arriba: cuadrícula → sombreados (dicfondo) → barras y áreas (dicbarras) → ejes, cero y umbrales
% (axis lines) → líneas (main) → puntos (dicpuntos) → rótulos (axis
% descriptions). No basta el orden del código: pgfplots dibuja cuadrícula,
% ejes y marcas en sus propias capas (y area style pone los ejes encima), y
% dentro de un axis los \draw, \fill y \node no admiten capa (van con los
% datos de main); por eso todo lo que se dibuja dentro del axis es un \addplot.
% Uso: \begin{axis}[…, dicapas] (al final de las opciones);
%   barras y áreas: \addplot[on layer=dicbarras, …];
%   sombreados:     \dicsombra[estilo]{x1}{y1}{x2}{y2} (rectángulo, en dicfondo,
%                   bajo las barras y sobre la cuadrícula);
%   cero y umbrales: \dicraya[estilo]{(x1,y1) (x2,y2)} (en axis lines);
%   líneas: \addplot normal (main); puntos: \addplot[on layer=dicpuntos, …]
%   (las marcas de cualquier serie van solas a dicpuntos);
%   rótulos: \dicrotulo[opciones del nodo]{(x,y)}{texto} (en axis descriptions),
%   siempre sin fondo (nada de fill=white: si chocan con una línea, se desplazan;
%   CRITERIOS §7, regla 4).
% Los \capa{n} del vídeo no cambian: siguen envolviendo instrucciones completas.
\pgfdeclarelayer{dicfondo}
\pgfdeclarelayer{dicbarras}
\pgfdeclarelayer{dicpuntos}
\pgfplotsset{
  % (el segundo argumento reparte cuadrícula, ejes y rótulos como el estándar)
  layers/dic/.define layer set={axis background, axis grid, dicfondo, dicbarras,
    axis ticks, axis lines, axis tick labels, pre main, main, dicpuntos,
    axis descriptions, axis foreground}{/pgfplots/layers/standard},
  % leyenda sin caja blanca: toma el fondo del gráfico, también si el gráfico va
  % dentro de un cuadro gris (CRITERIOS §7, regla 4, 26/10/2026)
  dicapas/.style={axis on top=false, set layers=dic, mark layer=dicpuntos,
    legend style={fill=none}},
  dicsuelto/.style={forget plot, stack plots=false, sharp plot, mark=none},
}
\newcommand{\dicraya}[2][]{%
  \addplot[on layer=axis lines, dicsuelto, diccero, #1] coordinates {#2};}
\newcommand{\dicsombra}[5][]{%
  \addplot[on layer=dicfondo, dicsuelto, draw=none, #1]
    coordinates {(#2,#3) (#4,#3) (#4,#5) (#2,#5)};}
\newcommand{\dicrotulo}[3][]{%
  \addplot[on layer=axis descriptions, dicsuelto, draw=none] coordinates {#2}
    node[pos=0, #1] {#3};}

%   \begin{document}
%   \begin{tikzpicture}[tcee]
%     \ejes{Q}{P}{6}{5}
%     \draw[curva1] (0.5,4.5) -- (5,0.8) node[etiqueta,right] {$D$};
%     \capa{2}{\draw[curva2] ...;}   % en el vídeo aparece en el paso 2
%   \end{tikzpicture}
%   \end{document}
%
% Paleta: la de la web (morado, dorado, salvia) más tres colores de apoyo
% distinguibles en blanco y negro por el trazo.
% ==========================================================================
\usepackage[T1]{fontenc}
\usepackage{newpxtext,newpxmath}
\usepackage{xcolor}
\usepackage{tikz,pgfplots}
\pgfplotsset{compat=1.18}
\usetikzlibrary{arrows.meta,calc,positioning,patterns,decorations.pathreplacing,intersections,shapes.misc,backgrounds}

\definecolor{tceemorado}{HTML}{5F2987}
\definecolor{tceedorado}{HTML}{B8860B}
\definecolor{tceeverde}{HTML}{3C7D3A}
\definecolor{tceeazul}{HTML}{1F5F99}
\definecolor{tceerojo}{HTML}{B22222}
\definecolor{tceenaranja}{HTML}{E07B28}
\definecolor{tceegris}{HTML}{767171}
\definecolor{tceesalvia}{HTML}{E2EFD9}

% Capas para el vídeo: \capa{n}{…} dibuja su contenido solo a partir del paso n.
% En el PDF, la web y Word se ve todo; generar-video.py compila el gráfico una
% vez por paso con \def\tceepaso{k} para que cada elemento aparezca al narrarlo.
% Vectores con flecha encima, como en el texto (tcee.sty): \vect{x}, \vect{p}
\providecommand{\vect}[1]{\vec{#1}}
\providecommand{\tceepaso}{1000}
\newcommand{\capa}[2]{\ifnum#1>\tceepaso\relax\else#2\fi}
% Gráficos por pasos del mismo tamaño (propuesta del 3.A.29, 11/11/2026): con
% \def\tceefantasma{1} antes del \input (p. ej., x-paso1.tex:
% «\def\tceepaso{3}\def\tceefantasma{1}\input{x.tex}»), las capas posteriores
% al paso se dibujan invisibles y ocupan su sitio. Sin \tceefantasma (el vídeo),
% no cambia nada.
\ifdefined\tceefantasma
  \renewcommand{\capa}[2]{\ifnum#1>\tceepaso\relax\begin{scope}[opacity=0]#2\end{scope}\else#2\fi}
\fi

\tikzset{
  tcee/.style={line cap=round,line join=round,font=\small,>={Stealth[length=2.2mm]}},
  eje/.style={->,thick,black},
  curva1/.style={very thick,tceemorado},
  curva2/.style={very thick,tceeazul},
  curva3/.style={very thick,tceeverde},
  curva4/.style={very thick,tceedorado},
  curvanueva/.style={very thick,tceerojo},
  desplazada/.style={very thick,dashed},
  guia/.style={thin,dashed,tceegris},
  punto/.style={circle,fill=black,inner sep=1.3pt},
  etiqueta/.style={font=\small,inner sep=2pt},
  flecha/.style={->,thick,tceerojo},
  area/.style={fill=tceemorado,fill opacity=0.18,draw=none},
  area2/.style={fill=tceedorado,fill opacity=0.22,draw=none},
}

% \ejes{etiqueta x}{etiqueta y}{largo x}{alto y}
\newcommand{\ejes}[4]{%
  \draw[eje] (0,0) -- (#3,0) node[below left=2pt and 0pt] {#1};
  \draw[eje] (0,0) -- (0,#4) node[left] {#2};
  \node[below left] at (0,0) {$0$};}
% \proyeccion{(x,y)}{etq x}{etq y}: líneas guía a los ejes con etiquetas
\newcommand{\proyeccion}[3]{%
  \draw[guia] let \p1=#1 in (\x1,0) node[below] {#2} |- (0,\y1) node[left] {#3};}

% --------------------------------------------------------------------------
% Paleta «dictamen» (dic…): gráficos de datos de la Guía del dictamen (1D01) y
% de cualquier gráfico de coyuntura que quiera los mismos colores. Un concepto,
% un color en todos los gráficos (CRITERIOS §7, «tres reglas», 24/10/2026);
% tabla completa en fuentes/1D01/_trabajo/paleta-1D01.md. No cambia nada de
% lo anterior: son colores y estilos nuevos.
% Actividad: el PIB (y sus variantes: PIB por persona, PIB mundial), morado.
\definecolor{dicpib}{HTML}{5F2987}
% Demanda interna, gama cálida: consumo (naranjas: privado oscuro, público
% claro), inversión (verde) y existencias (gris con trama, aparte).
\definecolor{dicconspriv}{HTML}{C76A1F}
\definecolor{dicconspub}{HTML}{F5C690}
\definecolor{dicfbcf}{HTML}{2A6B3C}
\definecolor{dicexistencias}{HTML}{B4B4B4}
% Demanda externa y comercio, gama fría (azul claro, que se separa del morado)
\definecolor{dicexterna}{HTML}{6AA6DA}
% Zonas: zona del euro (azul marino), España (rojo), resto de países (grises)
\definecolor{dicze}{HTML}{17406D}
\definecolor{dicespana}{HTML}{C8382E}
\definecolor{dicotro}{HTML}{8A8A8A}
\definecolor{dicotroclaro}{HTML}{B5B5B5}
\definecolor{dicotromasclaro}{HTML}{DADADA}
% Precios: inflación general (negro, continuo) y subyacente (negro,
% discontinuo); componentes: energía (bermellón), alimentos (amarillo) y los
% de la subyacente en azules (bienes industriales claro, servicios oscuro).
\definecolor{dicinflacion}{HTML}{1A1A1A}
\definecolor{dicenergia}{HTML}{D55E00}
\definecolor{dicalimentos}{HTML}{F2D45C}
\definecolor{dicbine}{HTML}{7FB2E0}
\definecolor{dicservicios}{HTML}{1F4E89}
% Trabajo: empleo (verde azulado), horas (el mismo, claro), tasa de actividad
% (verde azulado oscuro), paro (teja); productividad (dorado de la web).
\definecolor{dicempleo}{HTML}{1B7F6B}
\definecolor{dichoras}{HTML}{86CBB8}
\definecolor{dicactividad}{HTML}{0F5247}
\definecolor{dicparo}{HTML}{A23B2A}
\definecolor{dicproductividad}{HTML}{B8860B}
% Hogares: renta disponible (pizarra) y tasa de ahorro (negro); el consumo es
% dicconspriv.
\definecolor{dicrenta}{HTML}{4A5A6A}
\definecolor{dicahorro}{HTML}{1A1A1A}
% Política monetaria: tipos del BCE en azules (depósito, el de la zona del
% euro; operaciones principales, medio; facilidad marginal, claro); balance:
% préstamos a los bancos (dorado claro), resto del activo (gris claro).
\definecolor{dicdeposito}{HTML}{17406D}
\definecolor{dicopf}{HTML}{4F86C6}
\definecolor{dicfmc}{HTML}{9EC5EA}
\definecolor{dicprestamos}{HTML}{D9A520}
\definecolor{dicresto}{HTML}{CFCFCF}
% Ciclo: series adelantadas, coincidentes (el PIB) y retrasadas
\definecolor{dicadelantada}{HTML}{D55E00}
\definecolor{dicretrasada}{HTML}{6AA6DA}
% Referencias: rejilla, previsión, sombreados
\definecolor{dicreferencia}{HTML}{767171}

\tikzset{
  % barra de un componente: \addplot[dicbarra=dicfbcf]
  dicbarra/.style={fill=#1, draw=#1!65!black, line width=0.3pt},
  % existencias: gris con trama, para que se distinga también en gris
  dicbarratrama/.style={fill=#1, draw=#1!65!black, line width=0.3pt,
    postaction={pattern=north east lines, pattern color=white}},
  % línea de una serie; con halo blanco cuando pasa sobre barras
  diclinea/.style={very thick, #1},
  diclineahalo/.style={very thick, #1, preaction={draw=white, line width=3.4pt, solid}},
  % puntos de una serie (capa aparte, después de la línea)
  dicpuntos/.style={only marks, mark=*, mark size=1.8pt,
    mark options={fill=white, draw=#1, line width=0.9pt, solid}},
  % línea del cero o de un umbral, después de las barras
  diccero/.style={thin, black},
  dicumbral/.style={thin, dashed, dicreferencia},
}
% Capas de los gráficos de datos (CRITERIOS §7, «tres reglas»), de abajo
% arriba: cuadrícula → sombreados (dicfondo) → barras y áreas (dicbarras) → ejes, cero y umbrales
% (axis lines) → líneas (main) → puntos (dicpuntos) → rótulos (axis
% descriptions). No basta el orden del código: pgfplots dibuja cuadrícula,
% ejes y marcas en sus propias capas (y area style pone los ejes encima), y
% dentro de un axis los \draw, \fill y \node no admiten capa (van con los
% datos de main); por eso todo lo que se dibuja dentro del axis es un \addplot.
% Uso: \begin{axis}[…, dicapas] (al final de las opciones);
%   barras y áreas: \addplot[on layer=dicbarras, …];
%   sombreados:     \dicsombra[estilo]{x1}{y1}{x2}{y2} (rectángulo, en dicfondo,
%                   bajo las barras y sobre la cuadrícula);
%   cero y umbrales: \dicraya[estilo]{(x1,y1) (x2,y2)} (en axis lines);
%   líneas: \addplot normal (main); puntos: \addplot[on layer=dicpuntos, …]
%   (las marcas de cualquier serie van solas a dicpuntos);
%   rótulos: \dicrotulo[opciones del nodo]{(x,y)}{texto} (en axis descriptions),
%   siempre sin fondo (nada de fill=white: si chocan con una línea, se desplazan;
%   CRITERIOS §7, regla 4).
% Los \capa{n} del vídeo no cambian: siguen envolviendo instrucciones completas.
\pgfdeclarelayer{dicfondo}
\pgfdeclarelayer{dicbarras}
\pgfdeclarelayer{dicpuntos}
\pgfplotsset{
  % (el segundo argumento reparte cuadrícula, ejes y rótulos como el estándar)
  layers/dic/.define layer set={axis background, axis grid, dicfondo, dicbarras,
    axis ticks, axis lines, axis tick labels, pre main, main, dicpuntos,
    axis descriptions, axis foreground}{/pgfplots/layers/standard},
  % leyenda sin caja blanca: toma el fondo del gráfico, también si el gráfico va
  % dentro de un cuadro gris (CRITERIOS §7, regla 4, 26/10/2026)
  dicapas/.style={axis on top=false, set layers=dic, mark layer=dicpuntos,
    legend style={fill=none}},
  dicsuelto/.style={forget plot, stack plots=false, sharp plot, mark=none},
}
\newcommand{\dicraya}[2][]{%
  \addplot[on layer=axis lines, dicsuelto, diccero, #1] coordinates {#2};}
\newcommand{\dicsombra}[5][]{%
  \addplot[on layer=dicfondo, dicsuelto, draw=none, #1]
    coordinates {(#2,#3) (#4,#3) (#4,#5) (#2,#5)};}
\newcommand{\dicrotulo}[3][]{%
  \addplot[on layer=axis descriptions, dicsuelto, draw=none] coordinates {#2}
    node[pos=0, #1] {#3};}

%   \begin{document}
%   \begin{tikzpicture}[tcee]
%     \ejes{Q}{P}{6}{5}
%     \draw[curva1] (0.5,4.5) -- (5,0.8) node[etiqueta,right] {$D$};
%     \capa{2}{\draw[curva2] ...;}   % en el vídeo aparece en el paso 2
%   \end{tikzpicture}
%   \end{document}
%
% Paleta: la de la web (morado, dorado, salvia) más tres colores de apoyo
% distinguibles en blanco y negro por el trazo.
% ==========================================================================
\usepackage[T1]{fontenc}
\usepackage{newpxtext,newpxmath}
\usepackage{xcolor}
\usepackage{tikz,pgfplots}
\pgfplotsset{compat=1.18}
\usetikzlibrary{arrows.meta,calc,positioning,patterns,decorations.pathreplacing,intersections,shapes.misc,backgrounds}

\definecolor{tceemorado}{HTML}{5F2987}
\definecolor{tceedorado}{HTML}{B8860B}
\definecolor{tceeverde}{HTML}{3C7D3A}
\definecolor{tceeazul}{HTML}{1F5F99}
\definecolor{tceerojo}{HTML}{B22222}
\definecolor{tceenaranja}{HTML}{E07B28}
\definecolor{tceegris}{HTML}{767171}
\definecolor{tceesalvia}{HTML}{E2EFD9}

% Capas para el vídeo: \capa{n}{…} dibuja su contenido solo a partir del paso n.
% En el PDF, la web y Word se ve todo; generar-video.py compila el gráfico una
% vez por paso con \def\tceepaso{k} para que cada elemento aparezca al narrarlo.
% Vectores con flecha encima, como en el texto (tcee.sty): \vect{x}, \vect{p}
\providecommand{\vect}[1]{\vec{#1}}
\providecommand{\tceepaso}{1000}
\newcommand{\capa}[2]{\ifnum#1>\tceepaso\relax\else#2\fi}
% Gráficos por pasos del mismo tamaño (propuesta del 3.A.29, 11/11/2026): con
% \def\tceefantasma{1} antes del \input (p. ej., x-paso1.tex:
% «\def\tceepaso{3}\def\tceefantasma{1}\input{x.tex}»), las capas posteriores
% al paso se dibujan invisibles y ocupan su sitio. Sin \tceefantasma (el vídeo),
% no cambia nada.
\ifdefined\tceefantasma
  \renewcommand{\capa}[2]{\ifnum#1>\tceepaso\relax\begin{scope}[opacity=0]#2\end{scope}\else#2\fi}
\fi

\tikzset{
  tcee/.style={line cap=round,line join=round,font=\small,>={Stealth[length=2.2mm]}},
  eje/.style={->,thick,black},
  curva1/.style={very thick,tceemorado},
  curva2/.style={very thick,tceeazul},
  curva3/.style={very thick,tceeverde},
  curva4/.style={very thick,tceedorado},
  curvanueva/.style={very thick,tceerojo},
  desplazada/.style={very thick,dashed},
  guia/.style={thin,dashed,tceegris},
  punto/.style={circle,fill=black,inner sep=1.3pt},
  etiqueta/.style={font=\small,inner sep=2pt},
  flecha/.style={->,thick,tceerojo},
  area/.style={fill=tceemorado,fill opacity=0.18,draw=none},
  area2/.style={fill=tceedorado,fill opacity=0.22,draw=none},
}

% \ejes{etiqueta x}{etiqueta y}{largo x}{alto y}
\newcommand{\ejes}[4]{%
  \draw[eje] (0,0) -- (#3,0) node[below left=2pt and 0pt] {#1};
  \draw[eje] (0,0) -- (0,#4) node[left] {#2};
  \node[below left] at (0,0) {$0$};}
% \proyeccion{(x,y)}{etq x}{etq y}: líneas guía a los ejes con etiquetas
\newcommand{\proyeccion}[3]{%
  \draw[guia] let \p1=#1 in (\x1,0) node[below] {#2} |- (0,\y1) node[left] {#3};}

% --------------------------------------------------------------------------
% Paleta «dictamen» (dic…): gráficos de datos de la Guía del dictamen (1D01) y
% de cualquier gráfico de coyuntura que quiera los mismos colores. Un concepto,
% un color en todos los gráficos (CRITERIOS §7, «tres reglas», 24/10/2026);
% tabla completa en fuentes/1D01/_trabajo/paleta-1D01.md. No cambia nada de
% lo anterior: son colores y estilos nuevos.
% Actividad: el PIB (y sus variantes: PIB por persona, PIB mundial), morado.
\definecolor{dicpib}{HTML}{5F2987}
% Demanda interna, gama cálida: consumo (naranjas: privado oscuro, público
% claro), inversión (verde) y existencias (gris con trama, aparte).
\definecolor{dicconspriv}{HTML}{C76A1F}
\definecolor{dicconspub}{HTML}{F5C690}
\definecolor{dicfbcf}{HTML}{2A6B3C}
\definecolor{dicexistencias}{HTML}{B4B4B4}
% Demanda externa y comercio, gama fría (azul claro, que se separa del morado)
\definecolor{dicexterna}{HTML}{6AA6DA}
% Zonas: zona del euro (azul marino), España (rojo), resto de países (grises)
\definecolor{dicze}{HTML}{17406D}
\definecolor{dicespana}{HTML}{C8382E}
\definecolor{dicotro}{HTML}{8A8A8A}
\definecolor{dicotroclaro}{HTML}{B5B5B5}
\definecolor{dicotromasclaro}{HTML}{DADADA}
% Precios: inflación general (negro, continuo) y subyacente (negro,
% discontinuo); componentes: energía (bermellón), alimentos (amarillo) y los
% de la subyacente en azules (bienes industriales claro, servicios oscuro).
\definecolor{dicinflacion}{HTML}{1A1A1A}
\definecolor{dicenergia}{HTML}{D55E00}
\definecolor{dicalimentos}{HTML}{F2D45C}
\definecolor{dicbine}{HTML}{7FB2E0}
\definecolor{dicservicios}{HTML}{1F4E89}
% Trabajo: empleo (verde azulado), horas (el mismo, claro), tasa de actividad
% (verde azulado oscuro), paro (teja); productividad (dorado de la web).
\definecolor{dicempleo}{HTML}{1B7F6B}
\definecolor{dichoras}{HTML}{86CBB8}
\definecolor{dicactividad}{HTML}{0F5247}
\definecolor{dicparo}{HTML}{A23B2A}
\definecolor{dicproductividad}{HTML}{B8860B}
% Hogares: renta disponible (pizarra) y tasa de ahorro (negro); el consumo es
% dicconspriv.
\definecolor{dicrenta}{HTML}{4A5A6A}
\definecolor{dicahorro}{HTML}{1A1A1A}
% Política monetaria: tipos del BCE en azules (depósito, el de la zona del
% euro; operaciones principales, medio; facilidad marginal, claro); balance:
% préstamos a los bancos (dorado claro), resto del activo (gris claro).
\definecolor{dicdeposito}{HTML}{17406D}
\definecolor{dicopf}{HTML}{4F86C6}
\definecolor{dicfmc}{HTML}{9EC5EA}
\definecolor{dicprestamos}{HTML}{D9A520}
\definecolor{dicresto}{HTML}{CFCFCF}
% Ciclo: series adelantadas, coincidentes (el PIB) y retrasadas
\definecolor{dicadelantada}{HTML}{D55E00}
\definecolor{dicretrasada}{HTML}{6AA6DA}
% Referencias: rejilla, previsión, sombreados
\definecolor{dicreferencia}{HTML}{767171}

\tikzset{
  % barra de un componente: \addplot[dicbarra=dicfbcf]
  dicbarra/.style={fill=#1, draw=#1!65!black, line width=0.3pt},
  % existencias: gris con trama, para que se distinga también en gris
  dicbarratrama/.style={fill=#1, draw=#1!65!black, line width=0.3pt,
    postaction={pattern=north east lines, pattern color=white}},
  % línea de una serie; con halo blanco cuando pasa sobre barras
  diclinea/.style={very thick, #1},
  diclineahalo/.style={very thick, #1, preaction={draw=white, line width=3.4pt, solid}},
  % puntos de una serie (capa aparte, después de la línea)
  dicpuntos/.style={only marks, mark=*, mark size=1.8pt,
    mark options={fill=white, draw=#1, line width=0.9pt, solid}},
  % línea del cero o de un umbral, después de las barras
  diccero/.style={thin, black},
  dicumbral/.style={thin, dashed, dicreferencia},
}
% Capas de los gráficos de datos (CRITERIOS §7, «tres reglas»), de abajo
% arriba: cuadrícula → sombreados (dicfondo) → barras y áreas (dicbarras) → ejes, cero y umbrales
% (axis lines) → líneas (main) → puntos (dicpuntos) → rótulos (axis
% descriptions). No basta el orden del código: pgfplots dibuja cuadrícula,
% ejes y marcas en sus propias capas (y area style pone los ejes encima), y
% dentro de un axis los \draw, \fill y \node no admiten capa (van con los
% datos de main); por eso todo lo que se dibuja dentro del axis es un \addplot.
% Uso: \begin{axis}[…, dicapas] (al final de las opciones);
%   barras y áreas: \addplot[on layer=dicbarras, …];
%   sombreados:     \dicsombra[estilo]{x1}{y1}{x2}{y2} (rectángulo, en dicfondo,
%                   bajo las barras y sobre la cuadrícula);
%   cero y umbrales: \dicraya[estilo]{(x1,y1) (x2,y2)} (en axis lines);
%   líneas: \addplot normal (main); puntos: \addplot[on layer=dicpuntos, …]
%   (las marcas de cualquier serie van solas a dicpuntos);
%   rótulos: \dicrotulo[opciones del nodo]{(x,y)}{texto} (en axis descriptions),
%   siempre sin fondo (nada de fill=white: si chocan con una línea, se desplazan;
%   CRITERIOS §7, regla 4).
% Los \capa{n} del vídeo no cambian: siguen envolviendo instrucciones completas.
\pgfdeclarelayer{dicfondo}
\pgfdeclarelayer{dicbarras}
\pgfdeclarelayer{dicpuntos}
\pgfplotsset{
  % (el segundo argumento reparte cuadrícula, ejes y rótulos como el estándar)
  layers/dic/.define layer set={axis background, axis grid, dicfondo, dicbarras,
    axis ticks, axis lines, axis tick labels, pre main, main, dicpuntos,
    axis descriptions, axis foreground}{/pgfplots/layers/standard},
  % leyenda sin caja blanca: toma el fondo del gráfico, también si el gráfico va
  % dentro de un cuadro gris (CRITERIOS §7, regla 4, 26/10/2026)
  dicapas/.style={axis on top=false, set layers=dic, mark layer=dicpuntos,
    legend style={fill=none}},
  dicsuelto/.style={forget plot, stack plots=false, sharp plot, mark=none},
}
\newcommand{\dicraya}[2][]{%
  \addplot[on layer=axis lines, dicsuelto, diccero, #1] coordinates {#2};}
\newcommand{\dicsombra}[5][]{%
  \addplot[on layer=dicfondo, dicsuelto, draw=none, #1]
    coordinates {(#2,#3) (#4,#3) (#4,#5) (#2,#5)};}
\newcommand{\dicrotulo}[3][]{%
  \addplot[on layer=axis descriptions, dicsuelto, draw=none] coordinates {#2}
    node[pos=0, #1] {#3};}

\begin{document}
\newcommand{\solopaso}[2]{\ifnum\tceepaso>#1\relax\else#2\fi}
\begin{tikzpicture}[tcee,x=1cm,y=1cm,
    declare function={g(\t)=6*\t/(1+\t)-0.15*\t;}]
  % 1: ejes X e Y
  \capa{1}{\ejes{$X$}{$Y$}{7.2}{6.8}}
  \capa{1}{\node[etiqueta,font=\footnotesize,align=left,anchor=north west] at (7.25,0) {exportación de $A$\\importación de $B$};}
  \capa{1}{\node[etiqueta,font=\footnotesize,align=right,anchor=south east] at (-0.35,7.05) {importación de $A$\\exportación de $B$};}
  % 2: curvas de oferta neta de libre comercio, E_0 y la RRI de libre comercio
  \capa{2}{\draw[thick,tceemorado!55] plot[domain=0:6.6,samples=80] ({g(\x)},\x) node[etiqueta,above,text=tceemorado] {$\mathrm{ON}_A$};}
  \capa{2}{\draw[thick,tceeazul!55] plot[domain=0:6.9,samples=80] (\x,{g(\x)}) node[etiqueta,right,text=tceeazul] {$\mathrm{ON}_B$};}
  \capa{2}{\draw[guia,tceeverde] (0,0) -- (5.6,5.6) node[etiqueta,above right=-2pt,text=tceeverde] {$\mathrm{RRI}_0=\frac{p_X}{p_Y}$};}
  \capa{2}{\node[punto,label={[etiqueta]right=2pt:$E_0$}] (E0) at (4.217,4.217) {};}
  \capa{2}{\proyeccion{(E0)}{$X_0$}{$Y_0$}}
  % 2 (solo en el paso 1): CIC de A y de B tangentes a RRI_0 en E_0
  \capa{2}{\solopaso{2}{\draw[thick,tceedorado] plot[domain=1.2:-1.6,samples=40] ({4.217+\x},{4.217+\x+0.13*\x*\x}) node[etiqueta,left,font=\footnotesize,text=tceedorado] {$\mathrm{CIC}_A$};}}
  \capa{2}{\solopaso{2}{\draw[thick,tceeazul!70] plot[domain=-1.6:1.2,samples=40] ({4.217+\x+0.13*\x*\x},{4.217+\x}) node[etiqueta,right,font=\footnotesize,text=tceeazul] {$\mathrm{CIC}_B$};}}
  % 3: E_1, el arancel óptimo de A sobre ON_B (RRI_1 y CIC de A solo en el paso 2)
  \capa{3}{\node[punto,label={[etiqueta]above left=-1pt:$E_1$}] (E1) at (2.515,3.916) {};}
  \capa{3}{\solopaso{3}{\draw[guia,tceeverde] (0,0) -- (3.6,5.605) node[etiqueta,above,font=\footnotesize,text=tceeverde] {$\mathrm{RRI}_1$};}}
  \capa{3}{\solopaso{3}{\draw[thick,tceedorado] plot[domain=-1.6:1.4,samples=40] ({2.515+\x},{3.916+0.336*\x+0.13*\x*\x}) node[etiqueta,right,font=\footnotesize,text=tceedorado] {$\mathrm{CIC}_A'$};}}
  % 3: ON_A' contraída por el arancel de A
  \capa{3}{\draw[curvanueva] plot[domain=0:6.6,samples=80] ({0.6*g(\x)},\x) node[etiqueta,above,text=tceerojo] {$\mathrm{ON}_A'$};}
  \capa{3}{\draw[->,thin,tceegris] ({g(5.5)-0.12},5.5) -- ({0.6*g(5.5)+0.15},5.5);}
  % 4: ON_B' contraída por el arancel de B
  \capa{4}{\draw[curvanueva,desplazada] plot[domain=0:6.9,samples=80] (\x,{0.6*g(\x)}) node[etiqueta,right,text=tceerojo] {$\mathrm{ON}_B'$};}
  \capa{4}{\draw[->,thin,tceegris] (5.5,{g(5.5)-0.12}) -- (5.5,{0.6*g(5.5)+0.15});}
  % 5: equilibrio de Nash en aranceles y las 2 curvas de indiferencia comercial
  \capa{5}{\draw[thick,tceemorado] plot[domain=-1.5:1.2,samples=40] ({2.303+\x},{2.303+0.24*\x+0.13*\x*\x});}
  \capa{5}{\draw[thin,tceemorado] (3.56,2.81) -- (4.5,3.15) node[etiqueta,anchor=west,text=tceemorado] {$\mathrm{CIC}_A$};}
  \capa{5}{\draw[thick,tceeazul] plot[domain=-1.5:1.2,samples=40] ({2.303+0.24*\x+0.13*\x*\x},{2.303+\x});}
  \capa{5}{\draw[thin,tceeazul] (2.81,3.56) -- (3.05,4.5) node[etiqueta,anchor=south,text=tceeazul] {$\mathrm{CIC}_B$};}
  \capa{5}{\node[punto,label={[etiqueta]below right=1pt and 0pt:$E_N$}] (EN) at (2.303,2.303) {};}
  \capa{5}{\proyeccion{(EN)}{$X_N$}{$Y_N$}}
  \capa{5}{\draw[thin] (EN) ++(0.25,-0.5) -- (4.4,0.95) node[etiqueta,font=\footnotesize,align=left,anchor=west] {equilibrio de Nash en aranceles:\\menos comercio que en $E_0$};}
\end{tikzpicture}
\end{document}
